% Part: first-order-logic
% Chapter: sequent-calculus
% Section: soundness.tex

\documentclass[../../../include/open-logic-section]{subfiles}

\begin{document}

\iftag{FOL}
      {\olfileid[fr]{fol}{seq}{sou}}
      {\olfileid[fr]{pl}{seq}{sou}}

\olsection{Correction}

\begin{explain}
Un système de !!{derivation}, tel que le calcul des séquents, est
\emph{correct} s'il ne permet pas de !!{derive} des résultats qui
ne sont pas effectivement vrais. La correction constitue ainsi une
garantie pour un système de !!{derivation}. Selon la notion syntaxique
considérée, nous souhaitons notamment savoir que :
\begin{enumerate}
\item tout énoncé~$!A$ !!{derivable} sans hypothèses est
  \iftag{FOL}{valide}{une tautologie} ;
\item si !!a{sentence} est !!{derivable} à partir d'autres énoncés,
  il en est aussi une conséquence logique ;
\item si un ensemble d'!!{sentence}s est incohérent, il est insatisfaisable.
\end{enumerate}
Ces propriétés sont essentielles. Si l'une d'elles échoue, le système
de !!{derivation} est défectueux : il permet de !!{derive} trop de
résultats. Il est donc crucial d'établir sa correction.

Toutes ces notions syntaxiques sont définies par la !!{derivability}
de certains séquents dans le calcul des séquents. Pour démontrer
(1)--(3), nous devons donc établir une propriété sémantique des
séquents !!{derivable}s. Nous allons d'abord définir la \emph{validité}
d'un séquent, puis montrer que tout séquent !!{derivable} est valide.
Les propriétés (1)--(3) en découleront comme corollaires.
\end{explain}

\begin{defn}
\iftag{FOL}{!!^a{structure}~$\Struct M$}{!!^a{valuation}~$\pAssign{v}$}
\emph{satisfait} un séquent $\Gamma \Sequent \Delta$ si et seulement
si l'une au moins des conditions suivantes est vérifiée :
$\iftag{FOL}{\Sat/{M}}{\pSat/{v}}{!A}$ pour un certain $!A \in \Gamma$,
ou $\iftag{FOL}{\Sat{M}}{\pSat{v}}{!A}$ pour un certain $!A \in \Delta$.

Un séquent est \emph{valide} si et seulement si toute
\iftag{FOL}{!!{structure}~$\Struct M$}{!!{valuation}~$\pAssign{v}$}
le satisfait.
\end{defn}

\begin{thm}[Correction]
\ollabel{thm:sequent-soundness}
Si $\Log{LK}$ !!{derive}s $\Theta \Sequent \Xi$, alors
$\Theta \Sequent \Xi$ est valide.
\end{thm}

\begin{proof}
Soit $\pi$ une !!{derivation} de $\Theta \Sequent \Xi$.
Procédons par récurrence sur le nombre $n$ d'inférences de $\pi$.

Si ce nombre est nul, $\pi$ se réduit à un séquent initial.
Tout séquent initial $!A \Sequent !A$ est valide : pour toute
\iftag{FOL}{$\Struct M$}{$\pAssign{v}$}, on a soit
$\iftag{FOL}{\Sat/{M}}{\pSat/{v}}{!A}$, soit
$\iftag{FOL}{\Sat{M}}{\pSat{v}}{!A}$.
\iftag{prvTrue}{Le séquent initial $\quad \Sequent \ltrue$ est
également valide, puisque $\ltrue$ est toujours vrai.}{}
\iftag{prvFalse}{De même, $\lfalse \Sequent \quad$ est valide,
puisque $\lfalse$ est toujours faux.}{}
\footnote{Note de l'édition française : les cas des séquents initiaux
pour les constantes logiques, lorsqu'ils sont présents, ainsi que
les cas de l'échange et de la contraction ci-dessous, sont explicités ;
la source ne les traite pas séparément.}

Si le nombre d'inférences est strictement positif, distinguons les cas
selon la dernière inférence. Par hypothèse de récurrence, ses prémisses
sont valides, car la !!{derivation} de chacune d'elles contient
strictement moins de $n$ inférences.

Considérons d'abord les inférences à une seule prémisse. Un échange
ne modifie que l'ordre des énoncés ; une contraction ne modifie que
le nombre d'occurrences d'un même énoncé. Dans les deux cas, les
conditions de satisfaction de la prémisse et de la conclusion sont
équivalentes. La validité est donc préservée.
\begin{enumerate}
\item La dernière inférence est un affaiblissement. Alors
  $\Theta \Sequent \Xi$ est soit $!A, \Gamma \Sequent \Delta$
  si l'inférence est \LeftR{\Weakening}, soit
  $\Gamma \Sequent \Delta, !A$ si elle est \RightR{\Weakening}.
  La !!{derivation} se termine donc par l'une des formes suivantes :
  \begin{prooftree}
    \AxiomC{}
    \Deduce$\Gamma \fCenter \Delta$
    \RightLabel{\LeftR{\Weakening}}
    \UnaryInf$!A, \Gamma \fCenter \Delta$
    \DisplayProof\qquad\bottomAlignProof
    \AxiomC{}
    \Deduce$\Gamma \fCenter \Delta$
    \RightLabel{\RightR{\Weakening}}
    \UnaryInf$\Gamma \fCenter \Delta, !A$
  \end{prooftree}
  Par hypothèse de récurrence, $\Gamma \Sequent \Delta$ est valide.
  Pour toute \iftag{FOL}{!!{structure}~$\Struct{M}$}{!!{valuation}~$\pAssign{v}$},
  il existe donc soit $!C \in \Gamma$ tel que
  $\iftag{FOL}{\Sat/{M}}{\pSat/{v}}{!C}$, soit $!C \in \Delta$
  tel que $\iftag{FOL}{\Sat{M}}{\pSat{v}}{!C}$.

  Dans le premier cas, $!C \in \Theta$ aussi : l'affaiblissement
  à gauche ajoute un énoncé à $\Gamma$, tandis que l'affaiblissement
  à droite conserve cet antécédent. Dans le second cas, $!C \in \Xi$
  aussi : l'affaiblissement à droite ajoute un énoncé à $\Delta$,
  tandis que l'affaiblissement à gauche conserve ce conséquent.
  Nous avons donc soit $\iftag{FOL}{\Sat/{M}}{\pSat/{v}}{!C}$
  pour un certain $!C \in \Theta$, soit
  $\iftag{FOL}{\Sat{M}}{\pSat{v}}{!C}$ pour un certain $!C \in \Xi$.
  Ainsi, $\Theta \Sequent \Xi$ est valide.\footnote{Note de l'édition
  française : l'égalité $\Theta = !A, \Gamma$ donnée par la source
  ne vaut que pour l'affaiblissement à gauche. L'argument est explicité
  pour les deux règles.}
\item La dernière inférence est \LeftR{\lnot}. Sa prémisse est
  $\Gamma \Sequent \Delta, !A$ et sa conclusion est
  $\lnot !A, \Gamma \Sequent \Delta$. La !!{derivation} se termine par :
  \begin{prooftree}
    \AxiomC{}
    \Deduce$\Gamma \fCenter \Delta, !A$
    \RightLabel{\LeftR{\lnot}}
    \UnaryInf$\lnot !A, \Gamma \fCenter \Delta$
  \end{prooftree}
  On a $\Theta = \lnot !A, \Gamma$ et $\Xi = \Delta$.

  Par hypothèse de récurrence, $\Gamma \Sequent \Delta, !A$ est
  valide. Pour toute \iftag{FOL}{$\Struct{M}$}{$\pAssign{v}$},
  l'un au moins des cas suivants se présente : (a) il existe
  $!C \in \Gamma$ tel que $\iftag{FOL}{\Sat/{M}}{\pSat/{v}}{!C}$ ;
  (b) il existe $!C \in \Delta$ tel que
  $\iftag{FOL}{\Sat{M}}{\pSat{v}}{!C}$ ; (c)
  $\iftag{FOL}{\Sat{M}}{\pSat{v}}{!A}$.
  Montrons que $\Theta \Sequent \Xi$ est aussi valide.
  Soit \iftag{FOL}{une !!{structure}~$\Struct{M}$}{une !!{valuation}~$\pAssign{v}$}.
  Dans le cas (a), il existe $!C \in \Gamma$ tel que
  $\iftag{FOL}{\Sat/{M}}{\pSat/{v}}{!C}$, et $!C \in \Theta$ aussi.
  Dans le cas (b), il existe $!C \in \Delta$ tel que
  $\iftag{FOL}{\Sat{M}}{\pSat{v}}{!C}$, et $!C \in \Xi$ aussi.
  Enfin, dans le cas (c), $\iftag{FOL}{\Sat{M}}{\pSat{v}}{!A}$,
  donc $\iftag{FOL}{\Sat/{M}}{\pSat/{v}}{\lnot !A}$.
  Comme $\lnot !A \in \Theta$, il existe $!C \in \Theta$ tel que
  $\iftag{FOL}{\Sat/{M}}{\pSat/{v}}{!C}$.
  Ainsi, $\Theta \Sequent \Xi$ est valide.
\item La dernière inférence est \RightR{\lnot} : cas laissé en exercice.
\item La dernière inférence est \LeftR{\land}. Il y a deux variantes :
  on peut introduire $!A \land !B$ à gauche à partir de $!A$ ou
  de $!B$ placé à gauche dans la prémisse. Dans le premier cas,
  $\pi$ se termine par :
  \begin{prooftree}
    \AxiomC{}
    \Deduce$!A, \Gamma \fCenter \Delta$
    \RightLabel{\LeftR{\land}}
    \UnaryInf$!A \land !B, \Gamma \fCenter \Delta$
  \end{prooftree}
  On a $\Theta = !A \land !B, \Gamma$ et $\Xi = \Delta$.
  Considérons \iftag{FOL}{une !!{structure}~$\Struct M$}{une !!{valuation}~$\pAssign{v}$}.
  La validité de $!A, \Gamma \Sequent \Delta$, donnée par l'hypothèse
  de récurrence, entraîne l'un au moins des cas suivants :
  (a) $\iftag{FOL}{\Sat/{M}}{\pSat/{v}}{!A}$ ;
  (b) $\iftag{FOL}{\Sat/{M}}{\pSat/{v}}{!C}$ pour un certain
  $!C \in \Gamma$ ; (c) $\iftag{FOL}{\Sat{M}}{\pSat{v}}{!C}$
  pour un certain $!C \in \Delta$.
  Dans le cas (a), $\iftag{FOL}{\Sat/{M}}{\pSat/{v}}{!A \land !B}$.
  Il existe donc $!C \in \Theta$, à savoir $!A \land !B$, tel que
  $\iftag{FOL}{\Sat/{M}}{\pSat/{v}}{!C}$.
  Dans le cas (b), il existe $!C \in \Gamma$ tel que
  $\iftag{FOL}{\Sat/{M}}{\pSat/{v}}{!C}$, et $!C \in \Theta$ aussi.
  Dans le cas (c), il existe $!C \in \Delta$ tel que
  $\iftag{FOL}{\Sat{M}}{\pSat{v}}{!C}$, et $!C \in \Xi$ puisque
  $\Xi = \Delta$. Dans tous les cas,
  \iftag{FOL}{$\Struct M$}{$\pAssign{v}$} satisfait donc
  $!A \land !B, \Gamma \Sequent \Delta$.
  Puisque \iftag{FOL}{$\Struct{M}$}{$\pAssign{v}$} était arbitraire,
  $!A \land !B, \Gamma \Sequent \Delta$ est valide.\footnote{Note de
  l'édition française : la conclusion de cet argument rétablit
  la formule principale $!A \land !B$, omise à cet endroit dans la source.}
  Le cas où $!A \land !B$ est introduit à partir de $!B$ se traite
  de même, en remplaçant $!A$ par $!B$.
\item La dernière inférence est \RightR{\lor}. Il y a deux variantes :
  on peut introduire $!A \lor !B$ à droite à partir de $!A$ ou
  de $!B$ placé à droite dans la prémisse. Dans le premier cas,
  $\pi$ se termine par :
  \begin{prooftree}
    \AxiomC{}
    \Deduce$\Gamma \fCenter \Delta, !A$
    \RightLabel{\RightR{\lor}}
    \UnaryInf$\Gamma \fCenter \Delta, !A \lor !B$
  \end{prooftree}
  On a $\Theta = \Gamma$ et $\Xi = \Delta, !A \lor !B$.
  Considérons \iftag{FOL}{une !!{structure}~$\Struct{M}$}{une !!{valuation}~$\pAssign{v}$}.
  Comme $\Gamma \Sequent \Delta, !A$ est valide, l'un au moins
  des cas suivants se présente : (a) $\iftag{FOL}{\Sat{M}}{\pSat{v}}{!A}$ ;
  (b) $\iftag{FOL}{\Sat/{M}}{\pSat/{v}}{!C}$ pour un certain
  $!C \in \Gamma$ ; (c) $\iftag{FOL}{\Sat{M}}{\pSat{v}}{!C}$
  pour un certain $!C \in \Delta$.
  Dans le cas (a), $\iftag{FOL}{\Sat{M}}{\pSat{v}}{!A \lor !B}$.
  Dans le cas (b), il existe $!C \in \Gamma$ tel que
  $\iftag{FOL}{\Sat/{M}}{\pSat/{v}}{!C}$.
  Dans le cas (c), il existe $!C \in \Delta$ tel que
  $\iftag{FOL}{\Sat{M}}{\pSat{v}}{!C}$.
  Dans tous les cas, \iftag{FOL}{$\Struct{M}$}{$\pAssign{v}$}
  satisfait $\Gamma \Sequent \Delta, !A \lor !B$, c'est-à-dire
  $\Theta \Sequent \Xi$. Puisque
  \iftag{FOL}{$\Struct{M}$}{$\pAssign{v}$} était arbitraire,
  $\Theta \Sequent \Xi$ est valide. Le cas où $!A \lor !B$
  est introduit à partir de $!B$ se traite de même, en remplaçant
  $!A$ par $!B$.
\item La dernière inférence est \RightR{\lif}. Alors $\pi$ se termine par :
  \begin{prooftree}
    \AxiomC{}
    \Deduce$!A, \Gamma \fCenter \Delta, !B$
    \RightLabel{\RightR{\lif}}
    \UnaryInf$\Gamma \fCenter \Delta, !A \lif !B$
  \end{prooftree}
  Par hypothèse de récurrence, la prémisse est valide ; montrons que
  la conclusion l'est aussi. Soit
  \iftag{FOL}{$\Struct{M}$}{$\pAssign{v}$} arbitraire.
  Comme $!A, \Gamma \Sequent \Delta, !B$ est valide, l'un au moins
  des cas suivants se présente : (a) $\iftag{FOL}{\Sat/{M}}{\pSat/{v}}{!A}$ ;
  (b) $\iftag{FOL}{\Sat{M}}{\pSat{v}}{!B}$ ;
  (c) $\iftag{FOL}{\Sat/{M}}{\pSat/{v}}{!C}$ pour un certain
  $!C \in \Gamma$ ; (d) $\iftag{FOL}{\Sat{M}}{\pSat{v}}{!C}$
  pour un certain $!C \in \Delta$.
  Dans chacun des cas (a) et (b),
  $\iftag{FOL}{\Sat{M}}{\pSat{v}}{!A \lif !B}$ ; il existe donc
  $!C \in \Delta, !A \lif !B$ tel que
  $\iftag{FOL}{\Sat{M}}{\pSat{v}}{!C}$.
  Dans le cas (c), il existe $!C \in \Gamma$ tel que
  $\iftag{FOL}{\Sat/{M}}{\pSat/{v}}{!C}$.
  Dans le cas (d), il existe $!C \in \Delta$ tel que
  $\iftag{FOL}{\Sat{M}}{\pSat{v}}{!C}$.
  Dans tous les cas, \iftag{FOL}{$\Struct{M}$}{$\pAssign{v}$}
  satisfait $\Gamma \Sequent \Delta, !A \lif !B$.
  Puisque \iftag{FOL}{$\Struct{M}$}{$\pAssign{v}$} était arbitraire,
  $\Gamma \Sequent \Delta, !A \lif !B$ est valide.
  \iftag{FOL}{%
\item La dernière inférence est \LeftR{\lforall}. Il existe alors
  une !!{formula}~$!A(x)$ et un terme clos~$t$ tels que $\pi$
  se termine par :
\begin{prooftree}
    \AxiomC{}
    \Deduce$!A(t), \Gamma \fCenter \Delta$
    \RightLabel{\LeftR{\lforall}}
    \UnaryInf$\lforall[x][!A(x)], \Gamma \fCenter \Delta$
  \end{prooftree}
  Montrons que la conclusion $\lforall[x][!A(x)], \Gamma \Sequent \Delta$
  est valide. Considérons une !!{structure}~$\Struct M$.
  Puisque la prémisse $!A(t), \Gamma \Sequent \Delta$ est valide,
  l'un au moins des cas suivants se présente : (a) $\Sat/{M}{!A(t)}$ ;
  (b) $\Sat/{M}{!C}$ pour un certain $!C \in \Gamma$ ;
  (c) $\Sat{M}{!C}$ pour un certain $!C \in \Delta$.
  Dans le cas (a), d'après la \olref[syn][sem]{prop:quant-terms},
  si $\Sat{M}{\lforall[x][!A(x)]}$, alors $\Sat{M}{!A(t)}$.
  Comme $\Sat/{M}{!A(t)}$, on a donc
  $\Sat/{M}{\lforall[x][!A(x)]}$.
  Dans les cas (b) et (c), $\Struct{M}$ satisfait également
  $\lforall[x][!A(x)], \Gamma \Sequent \Delta$.
  Puisque $\Struct M$ était arbitraire,
  $\lforall[x][!A(x)], \Gamma \Sequent \Delta$ est valide.
\item La dernière inférence est \RightR{\lexists} : cas laissé en exercice.
\item La dernière inférence est \RightR{\lforall}. Il existe alors
  une !!{formula}~$!A(x)$ et une !!{constant}~$a$ telles que $\pi$
  se termine par :
  \begin{prooftree}
    \AxiomC{}
    \Deduce$\Gamma \fCenter \Delta, !A(a)$
    \RightLabel{\RightR{\lforall}}
    \UnaryInf$\Gamma \fCenter \Delta, \lforall[x][!A(x)]$
  \end{prooftree}
  La condition sur l'eigenvariable est satisfaite : $a$ ne figure ni
  dans $!A(x)$, ni dans $\Gamma$, ni dans $\Delta$.
  Par hypothèse de récurrence, la prémisse de la dernière inférence
  est valide. Montrons que la conclusion l'est aussi : pour toute
  !!{structure}~$\Struct M$, l'un au moins des cas suivants doit
  se présenter : (a) $\Sat{M}{\lforall[x][!A(x)]}$ ;
  (b) $\Sat/{M}{!C}$ pour un certain $!C \in \Gamma$ ;
  (c) $\Sat{M}{!C}$ pour un certain $!C \in \Delta$.

  Soit $\Struct{M}$ une !!{structure} arbitraire. Si (b) ou (c)
  est vérifié, nous avons le résultat voulu. Supposons donc que
  ni l'un ni l'autre ne le soit : pour tout $!C \in \Gamma$,
  $\Sat{M}{!C}$, et pour tout $!C \in \Delta$, $\Sat/{M}{!C}$.
  Il reste à établir (a), c'est-à-dire
  $\Sat{M}{\lforall[x][!A(x)]}$.
  D'après la \olref[syn][ass]{prop:sat-quant}, il suffit de montrer
  que $\Sat{M}{!A(x)}[s]$ pour toute valuation des variables~$s$.
  Soit donc $s$ une valuation des variables arbitraire.
  Considérons la structure $\Struct{M'}$ qui coïncide avec
  $\Struct{M}$, sauf que $\Assign{a}{M'} = s(x)$.
  D'après le \olref[syn][ext]{cor:extensionality-sent}, pour tout
  $!C \in \Gamma$, on a $\Sat{M'}{!C}$ puisque $a$ ne figure
  pas dans $\Gamma$ ; de même, pour tout $!C \in \Delta$,
  on a $\Sat/{M'}{!C}$.
  La prémisse étant valide, on obtient $\Sat{M'}{!A(a)}$.
  D'après la \olref[syn][ass]{prop:sentence-sat-true},
  $\Sat{M'}{!A(a)}[s]$, car $!A(a)$ est un énoncé.
  Or $\varAssign{s}{s}{x}$ et $s(x) = \Value{a}{M'}[s]$,
  par la définition même de $\Struct{M'}$.
  La \olref[syn][ext]{prop:ext-formulas} donne donc
  $\Sat{M'}{!A(x)}[s]$.
  Comme $a$ ne figure pas dans $!A(x)$, la
  \olref[syn][ext]{prop:extensionality} entraîne
  $\Sat{M}{!A(x)}[s]$. Puisque $s$ était arbitraire, nous avons
  montré $\Sat{M}{!A(x)}[s]$ pour toute valuation des variables.
\item La dernière inférence est \LeftR{\lexists} : cas laissé en exercice.
}{}
\end{enumerate}
Considérons maintenant les inférences à deux prémisses.
\begin{enumerate}
\item La dernière inférence est une coupure. Alors $\pi$ se termine par :
  \begin{prooftree}
    \AxiomC{}
    \Deduce$\Gamma \fCenter \Delta, !A$
    \AxiomC{}
    \Deduce$!A, \Pi \fCenter \Lambda$
    \RightLabel{\Cut}
    \BinaryInf$\Gamma, \Pi \fCenter \Delta, \Lambda$
  \end{prooftree}
  Soit \iftag{FOL}{une !!{structure}~$\Struct{M}$}{une !!{valuation}~$\pAssign{v}$}.
  Par hypothèse de récurrence, les deux prémisses sont valides ;
  \iftag{FOL}{$\Struct{M}$}{$\pAssign{v}$} les satisfait donc toutes deux.
  Distinguons les cas (a) $\iftag{FOL}{\Sat/{M}}{\pSat/{v}}{!A}$
  et (b) $\iftag{FOL}{\Sat{M}}{\pSat{v}}{!A}$.
  Dans le cas (a), pour satisfaire la prémisse de gauche,
  \iftag{FOL}{$\Struct{M}$}{$\pAssign{v}$} doit satisfaire
  $\Gamma \Sequent \Delta$, et satisfait alors aussi la conclusion.
  Dans le cas (b), pour satisfaire la prémisse de droite,
  \iftag{FOL}{$\Struct{M}$}{$\pAssign{v}$} doit satisfaire
  $\Pi \Sequent \Lambda$.\footnote{Note de l'édition française :
  le signe de différence ensembliste de la source est remplacé ici
  par le signe de séquent, conformément au raisonnement.}
  Là encore, \iftag{FOL}{$\Struct{M}$}{$\pAssign{v}$} satisfait la conclusion.
\item La dernière inférence est \RightR{\land}. Alors $\pi$ se termine par :
  \begin{prooftree}
    \AxiomC{}
    \Deduce$\Gamma \fCenter \Delta, !A$
    \AxiomC{}
    \Deduce$\Gamma \fCenter \Delta, !B$
    \RightLabel{\RightR{\land}}
    \BinaryInf$\Gamma \fCenter \Delta, !A \land !B$
  \end{prooftree}
  Considérons \iftag{FOL}{une !!{structure}~$\Struct M$}{une !!{valuation}~$\pAssign{v}$}.
  Si \iftag{FOL}{$\Struct{M}$}{$\pAssign{v}$} satisfait
  $\Gamma \Sequent \Delta$, le résultat est acquis.
  Supposons donc que ce ne soit pas le cas. Puisque
  $\Gamma \fCenter \Delta, !A$ est valide par hypothèse de récurrence,
  $\iftag{FOL}{\Sat{M}}{\pSat{v}}{!A}$.
  De même, la validité de $\Gamma \Sequent \Delta, !B$ entraîne
  $\iftag{FOL}{\Sat{M}}{\pSat{v}}{!B}$.
  Il en résulte $\iftag{FOL}{\Sat{M}}{\pSat{v}}{!A \land !B}$.
\item La dernière inférence est \LeftR{\lor} : cas laissé en exercice.
\item La dernière inférence est \LeftR{\lif}. Alors $\pi$ se termine par :
  \begin{prooftree}
    \AxiomC{}
    \Deduce$\Gamma \fCenter \Delta, !A$
    \AxiomC{}
    \Deduce$!B, \Pi \fCenter \Lambda$
    \RightLabel{\LeftR{\lif}}
    \BinaryInf$!A \lif !B, \Gamma, \Pi \fCenter \Delta, \Lambda$
  \end{prooftree}
  Considérons de nouveau
  \iftag{FOL}{une !!{structure}~$\Struct{M}$}{une !!{valuation}~$\pAssign{v}$}
  et supposons que \iftag{FOL}{$\Struct{M}$}{$\pAssign{v}$}
  ne satisfasse pas $\Gamma, \Pi \Sequent \Delta, \Lambda$.
  Il faut montrer $\iftag{FOL}{\Sat/{M}}{\pSat/{v}}{!A \lif !B}$.
  Si \iftag{FOL}{$\Struct{M}$}{$\pAssign{v}$} ne satisfait pas
  $\Gamma, \Pi \Sequent \Delta, \Lambda$, alors
  \iftag{FOL}{$\Struct{M}$}{$\pAssign{v}$} ne satisfait ni
  $\Gamma \Sequent \Delta$ ni $\Pi \Sequent \Lambda$.
  Puisque $\Gamma \Sequent \Delta, !A$ est valide,
  $\iftag{FOL}{\Sat{M}}{\pSat{v}}{!A}$.
  Puisque $!B, \Pi \Sequent \Lambda$ est valide,
  $\iftag{FOL}{\Sat/{M}}{\pSat/{v}}{!B}$.
  On a donc $\iftag{FOL}{\Sat/{M}}{\pSat/{v}}{!A \lif !B}$,
  comme voulu.
\end{enumerate}
\end{proof}

\tagprob{FOL}
\begin{prob}
Compléter la démonstration du \olref[fol][seq][sou]{thm:sequent-soundness}.
\end{prob}
\tagendprob

\tagprob{notFOL}
\begin{prob}
Compléter la démonstration du \olref[pl][seq][sou]{thm:sequent-soundness}.
\end{prob}
\tagendprob

\begin{cor}
\ollabel{cor:weak-soundness}
Si $\Proves !A$, alors $!A$ est \iftag{FOL}{valide}{une tautologie}.
\end{cor}

\begin{cor}
\ollabel{cor:entailment-soundness}
Si $\Gamma \Proves !A$, alors $\Gamma \Entails !A$.
\end{cor}

\begin{proof}
Si $\Gamma \Proves !A$, il existe une partie finie
$\Gamma_0 \subseteq \Gamma$ et une !!{derivation} de
$\Gamma_0 \Sequent !A$. D'après le \olref{thm:sequent-soundness},
toute \iftag{FOL}{!!{structure}~$\Struct{M}$}{!!{valuation}~$\pAssign{v}$}
rend faux un certain $!B \in \Gamma_0$ ou rend vrai $!A$.
Ainsi, si $\iftag{FOL}{\Sat{M}}{\pSat{v}}{\Gamma}$, alors
$\iftag{FOL}{\Sat{M}}{\pSat{v}}{!A}$ aussi.
\end{proof}

\begin{cor}
\ollabel{cor:consistency-soundness}
Si $\Gamma$ est satisfaisable, alors il est cohérent.
\end{cor}

\begin{proof}
Raisonnons par contraposition. Supposons $\Gamma$ incohérent.
Il existe alors une partie finie $\Gamma_0 \subseteq \Gamma$
et une !!{derivation} de $\Gamma_0 \Sequent \quad$.
D'après le \olref{thm:sequent-soundness}, ce séquent est valide.
Autrement dit, pour toute
\iftag{FOL}{!!{structure}~$\Struct{M}$}{!!{valuation}~$\pAssign{v}$},
il existe $!C \in \Gamma_0$ tel que
$\iftag{FOL}{\Sat/{M}}{\pSat/{v}}{!C}$.
Puisque $\Gamma_0 \subseteq \Gamma$, cet énoncé $!C$ appartient
aussi à $\Gamma$. Aucune \iftag{FOL}{$\Struct{M}$}{$\pAssign{v}$}
ne satisfait donc $\Gamma$, qui est ainsi insatisfaisable.
\end{proof}

\end{document}
